
Considere duas definições para o operador de produto interno entre vetores reais.

A primeira corresponde à definição usual da Álgebra Linear: para $x, y \in \mathbb{R}^n$,
\[
  \dotprod{x}{y}
  =
  \sum_{i=1}^{n} x_i y_i.
\]

A segunda é uma definição recursiva. Se $x$ e $y$ forem vetores vazios (de dimensão $0$), então
\[
  \dotprod{x}{y} = 0.
\]
Caso contrário, escrevendo
\[
  x
  =
  \begin{bmatrix}
    x_1 \\
    x_2
  \end{bmatrix},
  \qquad
  y
  =
  \begin{bmatrix}
    y_1 \\
    y_2
  \end{bmatrix},
\]
onde $x_1, y_1 \in \mathbb{R}$ e $x_2, y_2 \in \mathbb{R}^{n-1}$, define-se
\[
  \dotprod{x}{y}
  =
  x_1 y_1 + \dotprod{x_2}{y_2}.
\]

É natural esperar que ambas as definições sejam equivalentes. Entretanto, quando implementadas em uma linguagem imperativa como Julia, essa equivalência não fica evidente, pois a implementação da definição por somatório normalmente utiliza um laço de repetição:

\begin{code}
function dotprod1(xs, ys)
    s = 0
    for i in eachindex(xs)
        s = s + xs[i] * ys[i]
    end
    return s
end
\end{code}

Enquanto isso, a definição recursiva é implementada de forma bastante direta:

\begin{code}
function dotprod2(xs, ys)
    return xs[1] * ys[1] + dotprod1(xs[2:end], ys[2:end])
end
\end{code}

Embora ambos os programas calculem o produto interno, provar sua equivalência diretamente a partir dessas implementações é pouco natural, pois a estrutura do primeiro programa difere significativamente da definição matemática correspondente.

Em contraste, em Haskell as duas implementações refletem quase literalmente suas respectivas definições matemáticas:

\begin{code}
dotprod1 :: [Double] -> [Double] -> Double
dotprod1 xs ys = sum (zipWith (*) xs ys)
\end{code}

\begin{code}
dotprod2 :: [Double] -> [Double] -> Double
dotprod2 [] [] = 0
dotprod2 (x:xs) (y:ys) = x*y + dotprod1 xs ys
\end{code}

Nesse caso, uma prova por indução da equivalência entre os programas torna-se praticamente uma tradução da prova da equivalência entre as definições matemáticas.
